% ═══════════════════════════════════════════════════════════════════════════
%  DRAFT — expansion item #17: "A decision tree for practitioners"
%  Target: Chapter 6 (Conclusion and Operational Guidelines), insert after the
%  closing paragraph ("...artefacts of this specific sweep.") and before the
%  \newpage that precedes Appendix A (tex line ~2139).
%
%  STATUS: DRAFT for review. Not yet \input into dissertation.tex.
%  Evidence grades: A = repeated-CV / CI-backed or replicated across datasets;
%  B = single-split or r=5 with consistent replication; C = single observation.
%  Every branch cites the paper's own tables/figures (§4.1–§4.3, §5.2, Table
%  \ref{tab:stage_impact}, Figure \ref{fig:pareto}, Table \ref{tab:pareto}).
%
%  NOTE for the author: the reviewer's paraphrase "betti/PI top, landscape
%  weakest" does NOT match the paper's numbers — Persistence Landscapes is the
%  TOP vectorizer on ECG200 (75.32%, Table \ref{tab:stage_impact}) and
%  Persistence Entropy is the weakest (68.95%); Betti Curve + Silhouette top
%  ECG5000. The tree below uses the paper's actual numbers.
% ═══════════════════════════════════════════════════════════════════════════

\section{A Decision Tree for Practitioners}
\label{sec:decision_tree}

The conditionality established in this dissertation can be collapsed into a
small decision tree. Each branch names the evidence that supports it and the
grade of that evidence; grades follow the definitions above. The tree is
deliberately conservative --- where the evidence is thin, the branch says so
rather than pretending otherwise.

\begin{enumerate}
  \item \textbf{Modality.} The first question is the data modality.
    \begin{enumerate}
      \item \emph{Time series (after Takens embedding, $d{=}3$, $\tau{=}1$):}
        spend your design budget on the \textbf{vectorizer}, not the
        filtration. Vectorization is the dominant stage (ECG200 marginal range
        $6.39$pp, 95\% CI [6.13, 6.65]; ECG5000 $24.89$pp; \S4.1, Table
        \ref{tab:stage_impact}), while filtration contributes $0.69$pp
        (CI [0.57, 0.81]) on ECG200. \emph{Grade A} (r$=$25, Nadeau--Bengio
        corrected CIs exclude zero for the vectorizer only).
      \item \emph{Images:} if the task is \textbf{binary}, vectorization still
        leads (MNIST binary $3.22$pp vs.\ filtration $1.65$pp; r$=$5) but
        filtration is no longer negligible --- prefer the \textbf{cubical}
        filtration over Vietoris--Rips ($\sim$1.75pp best-of-family, 98.0\%
        vs.\ 96.25\%; \S4.2). \emph{Grade B} (r$=$5, stable across all five
        repetitions). If the task is \textbf{multi-class} ($\ge$10 classes),
        the ordering flips: filtration becomes the larger marginal stage
        (10-class MNIST: $4.53$pp vs.\ vectorizer $3.44$pp), though the swing
        (9--10pp within a filtration) is far below Conti et al.'s 18--94\%
        grid-search result (\S4.2). \emph{Grade B} (single sweep, 1000 samples).
      \item \emph{Point clouds:} if the classes differ in topology (genus,
        number of components) the pipeline separates them nearly perfectly and
        no stage dominates (sphere/torus $\ge$99.5\% at $\sigma{=}0$, mean
        99.99\%; \S4.2). If you must choose a stage to tune, choose the
        classifier (its $3.50$pp range is the only non-negligible one on
        saturated tasks). \emph{Grade A} for saturation; the matched-genus
        control shows TDA retains 95.83\% at $\sigma{=}0.30$ where norm
        features collapse to 48--58\% (\S4.2). \emph{Grade A}.
    \end{enumerate}

  \item \textbf{Budget: accuracy vs.\ runtime.} If wall-clock time is the
    constraint:
    \begin{enumerate}
      \item \emph{Large $n$ ($\ge$10$^3$ points):} Sparse Rips is the only
        filtration whose design point applies; Vietoris--Rips at $n{=}3000$ is
        infeasible ($O(n^3)$ 2-simplices). \emph{Grade B} (B5 large-n sweep;
        runtime row of the recommendation table).
      \item \emph{Small $n$:} any of Vietoris--Rips, weak Alpha, or Cubical;
        the 3--27\% wall-clock gap between them (\S4.3) rarely justifies
        accuracy risk. Avoid weak Alpha on quantized UCR series (giotto
        crashes; fragility is a first-class result, expansion \#11).
        \emph{Grade A} (wall-clock recorded per configuration in the results
        DBs).
      \item \emph{Classifier:} prefer Random Forest or logistic regression on
        TDA features; SVM-RBF collapses to the majority class (66.5\% on
        ECG200) on both TDA-only and concatenated features (\S5.2).
        \emph{Grade B}.
    \end{enumerate}

  \item \textbf{Noise level.} In the studied regime ($\sigma \le 0.30$
    additive spatial Gaussian noise, $n{=}100$), noise robustness is not a
    differentiator: mean accuracy across the 112 configurations at
    $\sigma{=}0.30$ is 99.85\% (min 98.5\%), and measured bottleneck distances
    (max 0.434) sit far below the corrected stability bound
    ($2\sigma\sqrt{2\ln n} \approx 1.82$) --- the bound is conservative, not
    tight (\S4.2). If your application lives in this regime, do not trade
    accuracy for noise-robust filtrations (e.g.\ DTM-weighting gained range
    but not headroom on ECG200; B1). \emph{Grade A}.

  \item \textbf{Which vectorizer?} Prefer \emph{distributional} vectorizers
    over scalar summaries:
    \begin{enumerate}
      \item Persistence Entropy --- the single-scalar-per-dimension
        representation --- is the weakest vectorizer on ECG200 (69.3\%
        single-split marginal; 68.95\% in the r$=$25 table) and catastrophic
        on ECG5000 (36.9\% vs.\ 61.5\% for Betti Curve). \emph{Grade A}.
      \item Top performers by dataset: Persistence Landscapes (ECG200,
        75.32\%), Betti Curve and Silhouette (ECG5000, 61.5\%/61.8\%),
        Betti Curve (10-class MNIST, 33.4\%). \emph{Grade B}.
      \item Persistence Statistics is a zero-parameter, cheap default: within
        0.3pp of Landscapes on ECG200 (74.8\%) and 2.1pp of the best MNIST
        configuration (95.9\%; \S4.1). \emph{Grade B}.
    \end{enumerate}

  \item \textbf{Raw features vs.\ TDA.} On the datasets studied, raw features
    beat TDA alone (ECG200 raw logistic 85.28\% vs.\ 83.0\% best TDA under
    r$=$25; MNIST raw-pixel logistic 99.65\% vs.\ 98.0\%; \S4.3). Do not use
    TDA as your \emph{only} feature source on such data. If you want the
    geometric signal, \textbf{concatenate}: $[\text{raw} \| \text{TDA}]$
    adds 0.35--1.72pp over the 25-repetition raw baselines (ECG200 87.0\%
    vs.\ 85.28\%; MNIST 100.0\% vs.\ 99.65\%), never hurts with logistic or
    Random Forest, and the benefit is confined to those two classifiers
    (\S5.2). \emph{Grade B} (single protocol, both datasets).

  \item \textbf{Is topology the target?} If the scientific question is about
    shape --- cycle structure, genus, component counts (dynamical systems,
    shape/geometry benchmarks) --- TDA is not a feature extractor to be
    benchmarked against raw baselines; it is the object of study. The
    matched-genus experiment is the proof of concept: TDA retains 95.83\%
    where norm/scale features fail (48--58\%). The topology-wins regime
    (expansion \#6) is designed to map this branch with the full stage
    decomposition; until it runs, treat this branch as \emph{Grade C}.
\end{enumerate}

\begin{table}[H]
  \centering
  \caption{Decision tree summary with evidence grades (expansion \#17 draft).}
  \label{tab:decision_tree}
  \begin{tabularx}{\textwidth}{@{}l l X c@{}}
    \toprule
    Branch & Decision & Evidence (paper) & Grade \\
    \midrule
    Time series & tune vectorizer & ECG200 6.39pp [6.13, 6.65] r$=$25; ECG5000 24.89pp (\S4.1, Table \ref{tab:stage_impact}) & A \\
    Image, binary & vectorizer first; cubical $>$ VR & MNIST 3.22 vs.\ 1.65pp r$=$5; 98.0 vs.\ 96.25\% (\S4.2) & B \\
    Image, multi-class & filtration first & 10-class MNIST 4.53 vs.\ 3.44pp (\S4.2) & B \\
    Point cloud & no stage dominates & $\ge$99.5\% all configs, $\sigma{=}0$; 99.85\% at $\sigma{=}0.30$ (\S4.2) & A \\
    Speed, large $n$ & Sparse Rips & B5 large-n sweep (design point) & B \\
    Speed, small $n$ & VR/weak Alpha/Cubical; no weak Alpha on quantized UCR & wall-clock in DBs (\S4.3); fragility finding & A \\
    Classifier & RF / logistic; not SVM-RBF on TDA features & 66.5\% majority-class collapse (\S5.2) & B \\
    Noise $\le 0.30$ & no robustness trade-off needed & 99.85\% mean; bottleneck 0.434 $\ll$ 1.82 (\S4.2) & A \\
    Vectorizer & distributional $>$ scalar; entropy last & entropy 68.95\%/36.9\%; landscape 75.32\% (Tables \ref{tab:stage_impact}, \S4.1--\S4.2) & A \\
    Raw available & concat $[\text{raw}\|\text{TDA}]$ with RF/logistic & +0.35--1.72pp over raw baselines (\S5.2) & B \\
    Topology is the target & TDA as object of study & matched-genus 95.83 vs.\ 48--58\% (\S4.2); \#6 pending & C \\
    \bottomrule
  \end{tabularx}
\end{table}

%% END OF DRAFT — insertion point: Chapter 6, before the \newpage at tex ~line 2139.
